% Résumé du mémoire.
% Abstract in French.
%
\chapter*{RÉSUMÉ}\thispagestyle{headings}
\addcontentsline{toc}{compteur}{RÉSUMÉ}

Les simulations tridimensionnelles et multicouches d’accrétion de givre \addText{sur les aéronefs} présentent d’importants défis numériques puisqu’elles nécessitent le couplage successif de modèles aérodynamiques, d’impact de gouttelettes, de thermodynamique et d’évolution géométrique, tout en mettant à jour de façon répétée une géométrie givrée de plus en plus complexe. En particulier, les méthodes basées sur le déplacement explicite de la surface et le remaillage répété de maillages conformes à la géométrie peuvent devenir sensibles à la qualité du maillage lorsque des structures 3D complexes se développent. L'objectif de cette thèse est d'améliorer la robustesse des simulations 3D et multicouches d'accrétion de givre par le développement et l'évaluation de méthodologies numériques complémentaires au sein de \textsc{CHAMPS}, un logiciel multiphysique développé à l’interne à Polytechnique Montréal.

Dans un premier temps, une méthodologie lagrangienne de suivi de particules est développée pour la prédiction de l'impact des gouttelettes. L'approche combine une intégration implicite des trajectoires, une détection accélérée des intersections, une injection adaptative des gouttelettes ainsi que des techniques d'interpolation à la surface afin d'améliorer l'efficacité numérique et la prédiction de la quantité d'eau collectée à la surface. Les effets associés aux grosses gouttelettes surfondues (SLD), notamment la déformation des gouttelettes, l'éclaboussement et le suivi des gouttelettes secondaires, sont intégrés aux formulations eulérienne et lagrangienne. Les prédictions obtenues montrent un bon accord avec les données de référence et démontrent une bonne cohérence entre les approches eulérienne et lagrangienne. La comparaison met également en évidence leur complémentarité : la formulation eulérienne offre un coût de calcul plus prévisible et une parallélisation efficace pour les grands problèmes 3D, tandis que la formulation lagrangienne donne directement accès aux trajectoires individuelles et se prête particulièrement bien à l'intégration de phénomènes physiques à l'échelle des gouttelettes.

Dans un deuxième temps, une méthodologie d’évolution géométrique par level set est développée pour l'accrétion de givre 3D et multicouche. L'interface givre--air est représentée implicitement par l'isocontour nul d'un champ de distance signée défini sur un maillage volumique auxiliaire. La méthodologie \removeText{advecte l’interface}\addText{utilise un schéma temporel semi-implicite (Rosenbrock-Wanner) pour advecter l’interface} selon une vitesse déterminée à partir des quantités de givrage, puis \addText{utilise la bibliothèque MMG pour }extraire l’interface afin d’obtenir une surface conforme à la géométrie résultante. Les cas de référence démontrent \addText{la précision du solveur et sa capacité à reproduire des structures 3D complexes, notamment les scallops observés expérimentalement,} \removeText{que l'approche permet de reproduire les formes de givre attendues} tout en offrant une meilleure robustesse géométrique que le déplacement explicite des nœuds de surface. Cet avantage devient particulièrement important lorsque la complexité géométrique augmente, puisque la représentation implicite évite l'enchevêtrement du maillage surfacique et réduit le recours à des corrections géométriques heuristiques au cours des couches successives d'accrétion.

Finalement, \removeText{le cadre numérique de givrage est étendu à une formulation par frontière immergée}\addText{les avancées réalisées en matière de suivi des gouttelettes et d'évolution géométrique par level set sont intégrées dans un cadre numérique de givrage reposant sur une approche hybride par frontière immergée}, dans laquelle les solutions aérodynamiques et d'impact de gouttelettes sont découplées d'un maillage volumique conforme à la géométrie. \addText{Les calculs thermodynamiques et l'évolution géométrique conservent toutefois une représentation surfacique conforme afin de préserver les avantages de la méthodologie d'évolution géométrique par level set.} Les cas de vérification et d'accrétion montrent un accord encourageant entre les prédictions obtenues avec les approches à maillage conforme et à frontière immergée pour l'impact des gouttelettes et la croissance du givre. Les résultats mettent également en évidence l'importance de la résolution locale du maillage ainsi que du traitement aérodynamique et thermique près de la paroi. Le raffinement adaptatif du maillage apparaît notamment comme un élément essentiel pour étendre efficacement la méthodologie par frontière immergée à des configurations 3D complexes.

Dans son ensemble, cette thèse établit un ensemble d'outils numériques complémentaires visant à améliorer la robustesse et la flexibilité des simulations 3D d'accrétion de givre. Les résultats montrent que l'évolution géométrique par level set constitue une base robuste pour les simulations multicouches, tandis que le suivi lagrangien des gouttelettes et les méthodes par frontière immergée offrent des capacités supplémentaires pour la modélisation détaillée des gouttelettes et le traitement de géométries complexes.